\section{Exploração dos conjuntos de dados}
\label{sec:datasets}

O dataset precisa, antes de tudo, \textbf{se encaixar no paralelismo de
dados}. Os critérios de escolha são:

\begin{enumerate}
  \item \textbf{Registros independentes:} cada linha pode ser processada
    sem olhar as outras, o que permite cortar os dados em pedaços de
    qualquer tamanho.
  \item \textbf{Resultado parcial combinável:} soma, contagem, mínimo e
    máximo se combinam diretamente; a média também, se cada worker devolver
    soma e contagem. Já a \emph{mediana} não pode ser obtida das medianas
    dos pedaços, um exemplo dos limites do modelo.
  \item \textbf{Volume suficiente:} cerca de 100\,MB para validar e de 1 a
    5\,GB para os experimentos; abaixo disso, o custo de distribuir supera
    o ganho.
  \item \textbf{Acesso aberto e reprodutível:} download público, com o
    script de obtenção versionado no repositório.
\end{enumerate}

\begin{table}[H]
  \centering
  \small
  \begin{tabularx}{\textwidth}{>{\raggedright}p{2.8cm}X>{\raggedright\arraybackslash}p{3cm}}
    \toprule
    \textbf{Dataset} & \textbf{Perguntas possíveis} & \textbf{Processamento} \\
    \midrule
    Microdados do ENEM (INEP)~\cite{enem} &
      Média de notas por estado, tipo de escola ou faixa de renda &
      Agregação por grupo \\
    NYC Taxi Trip Records~\cite{nyctaxi} &
      Corridas e receita por hora do dia ou por região &
      Agregação e filtragem \\
    \emph{Dumps} da Wikipédia &
      Palavras mais frequentes &
      Contagem (\emph{word count}) \\
    Dados sintéticos (\emph{seed} fixa) &
      Qualquer, com volume controlado &
      Qualquer \\
    \bottomrule
  \end{tabularx}
  \caption{Datasets candidatos.}
\end{table}

O grupo priorizou \textbf{primeiro a técnica e depois o negócio}: o núcleo
do trabalho (distribuição, agregação e medições) não depende do tema dos
dados, desde que o dataset atenda aos critérios acima. Por isso, o
desenvolvimento pode começar com dados sintéticos enquanto os candidatos
reais são analisados. O dataset ainda será decidido; a tendência é
combinar dados reais, que dão sentido à pergunta, com sintéticos, que
permitem variar o volume nos experimentos de escala. Os dados não serão
versionados no Git, apenas os scripts que os obtêm ou geram.
